% Generated draft: requires mathematical and source review.
\documentclass[11pt]{article}
\usepackage[margin=25mm]{geometry}
\usepackage{fontspec}
\setmainfont{Latin Modern Roman}
\usepackage{amsmath,amssymb,mathtools,unicode-math}
\setmathfont{Latin Modern Math}
\usepackage{enumitem,needspace,xurl,hyperref}
\setlength{\parindent}{0pt}
\setlength{\parskip}{5pt}
\setlength{\emergencystretch}{3em}
\newcommand{\N}{\mathbb{N}}
\newcommand{\Z}{\mathbb{Z}}
\newcommand{\Q}{\mathbb{Q}}
\newcommand{\R}{\mathbb{R}}
\newcommand{\C}{\mathbb{C}}
\newcommand{\F}{\mathbb{F}}
\newcommand{\A}{\mathbb{A}}
\newcommand{\PP}{\mathbb{P}}
\newcommand{\E}{\mathbb{E}}
\newcommand{\eps}{\varepsilon}
\newcommand{\dd}{\,\mathrm{d}}
\begin{document}
\section*{UnsolvedMath ID: 30001973}
\textbf{Source problem code:} catalogue entry
\subsection{Definitions and mathematical statement}
\textbf{Nature of the statement.} The supplied record contains a historical research question about nineteenth-century mathematical education. It supplies no formal mathematical definitions, mathematical domains, or quantified proposition. Its statement-verification note characterizes the question as a historical, pedagogical, or methodological discussion prompt rather than an unresolved mathematical proposition, and its statement status is recorded as \texttt{unrecoverable}.

\textbf{Preserved research question.} What mathematical practices were students encouraged to master in nineteenth-century mathematical education, and what methods, priorities, and evaluation criteria governed that training?

\textbf{Unresolved scope.} The record does not specify geographical scope, institutions, educational levels, student populations, documentary evidence, or operational meanings of mathematical practices, mastery, priorities, and evaluation criteria. Consequently, the question cannot be converted into a precise mathematical proposition from the supplied evidence. These qualifications remain unresolved; no formal hypotheses or quantifiers are added.
\subsection{Short English statement}
The question asks which mathematical practices students were encouraged to master in nineteenth-century education and how their training was organized and evaluated. The supplied record treats this as an interpretive historical research question. Its supplied open label does not establish an open mathematical conjecture.
\subsection{Sources}
\begin{itemize}
\item Local UnsolvedMath registry, \texttt{data/unsolvedmath/problems.json}, permanent ID \texttt{30001973}, title \emph{Practices of Nineteenth-Century Mathematical Education}, source code \texttt{catalogue entry}. The supplied record belongs to the Oberwolfach Reports Open Problems collection; no individual proposer is identified, and the recorded proposed year is 2012.
\item Supplied report citation: \emph{Explicit Versus Tacit Knowledge in Mathematics} (2012), Oberwolfach Reports, pp.~131--210, DOI 10.4171/owr/2012/04. \url{https://doi.org/10.4171/owr/2012/04}. This same link is supplied under literature sources as the ``original Oberwolfach report.'' The background describes Oberwolfach Reports as documenting mathematical workshops and their open-problem discussions; no specific page locating this question is supplied.
\item Supplied literature assessment dated exactly \texttt{2026-08-21}: status recorded as open; the extracted question concerns mathematical practices in nineteenth-century training and remains an interpretive historical research question. The original Oberwolfach report is listed as the source checked. This is supplied dated curation and literature triage, not independently verified here; its time-sensitive status and cited claims require independent verification. The supplied statement-verification note further characterizes the official history-of-mathematics workshop material as a historical, pedagogical, or methodological discussion prompt, promotes no separate clean statement, and retains the original record for provenance.
\end{itemize}
\end{document}
